\begin{ouverture}[grandeurs]{12}{Horaires et durées}{Partie 2 · Modéliser et mesurer}
{Instants · Intervalles · Base soixante · Heures décimales}
\objectif{Distinguer un horaire d’une durée et lire l’heure.}
\objectif{Calculer un début, une fin ou un temps écoulé.}
\objectif{Convertir des durées et relier heures décimales et heures-minutes.}
\prerequis{Lecture de l’heure ; jour, heure, minute, seconde.}
\end{ouverture}
\pagestyle{otMain}\essentiel\fontsize{11}{15}\selectfont
\phantomsection\label{cours-12-1}
\essSection{1}{Lire et représenter un horaire}
\begin{definition}Un horaire indique un instant : 14 h 32 min 5 s. Une durée mesure le temps écoulé entre deux instants : 1 h 28 min.\end{definition}
\begin{propriete}1 jour $=24$ h ; 1 h $=60$ min ; 1 min $=60$ s ; donc 1 h $=3\,600$ s. Une année compte 365 jours, ou 366 si elle est bissextile ; un siècle dure 100 ans, un millénaire $1\,000$ ans. Une demi-heure dure 30 min, un quart d’heure 15 min, trois quarts d’heure 45 min.\end{propriete}
\essSection{2}{Calculer un temps écoulé}
\begin{savoirfaire}[Avancer par étapes]De 9 h 47 à 11 h 15 : 13 min pour atteindre 10 h, puis 1 h pour atteindre 11 h, puis 15 min. Durée : $13+15=28$ min et 1 h, soit 1 h 28 min.\end{savoirfaire}
\label{fig-cours-12-1-cours-frise}\figureGeo{cours-frise}

\begin{exemple}Un film commence à 20 h 50 et dure 1 h 45 min. 10 min mènent à 21 h ; il reste 1 h 35 min : le film finit à 22 h 35.\end{exemple}
\clearpage
\phantomsection\label{cours-12-2}
\essSection{3}{Convertir des durées}
\begin{savoirfaire}[Diviser par 60]200 min : $200=60\times3+20$, soit 3 h 20 min. De même, $1\,000$ s : $1\,000=60\times16+40$, soit 16 min 40 s. Les durées se comptent par soixante, comme chez les Babyloniens il y a quatre mille ans.\end{savoirfaire}
\essSection{4}{Relier les écritures d’une durée}
\begin{propriete}Une durée peut s’écrire en heures décimales : $1{,}5$ h $=1$ h 30 min, car $0{,}5$ h est la moitié de 60 min ; $0{,}25$ h $=15$ min ; $0{,}1$ h $=6$ min. Inversement, 2 h 45 min $=2{,}75$ h.\end{propriete}
\label{fig-cours-12-2-cours-cadran}\figureGeo{cours-cadran}

\begin{attention}$1{,}30$ h ne vaut pas 1 h 30 min : $0{,}3$ h, c’est $0{,}3\times60=18$ min. Donc $1{,}3$ h $=1$ h 18 min.\end{attention}

\clearpage\rubrique{Découverte}
\begin{blocLivre}{ex-12-1}\begin{exercice}[1]{Lire l’heure}
Les horloges indiquent des horaires de l’après-midi. Les écrire en heures et minutes, de 0 h à 24 h.
\label{fig-ex-12-1}\figureGeo{dur-1}
\end{exercice}\end{blocLivre}
\begin{blocLivre}{ex-12-2}\begin{exercice}[1]{Unités de durée}
Compléter : 1 h $=\square$ s ; 3 jours $=\square$ h ; 2 min 15 s $=\square$ s ; 1 millénaire $=\square$ siècles.
\end{exercice}\end{blocLivre}
\begin{blocLivre}{ex-12-3}\begin{exercice}[1]{Fractions d’heure}
Combien de minutes dans : un quart d’heure ; trois quarts d’heure ; une demi-heure ; un tiers d’heure ?
\end{exercice}\end{blocLivre}
\begin{blocLivre}{ex-12-4}\begin{exercice}[1]{Calculer une durée}
Un cours de musique commence à 9 h 47 et se termine à 11 h 15. Combien de temps dure-t-il ?
\end{exercice}\end{blocLivre}
\begin{blocLivre}{ex-12-5}\begin{exercice}[1]{Heure d’arrivée}
Un car part à 14 h 35. Le trajet dure 2 h 50 min. À quelle heure arrive-t-il ?
\end{exercice}\end{blocLivre}

\rubrique{Entraînement}
\begin{blocLivre}{ex-12-6}\begin{exercice}[2]{Convertir}
Convertir en heures et minutes : 200 min ; 135 min. Convertir $1\,000$ s en minutes et secondes.
\end{exercice}\end{blocLivre}
\begin{blocLivre}{ex-12-7}\begin{exercice}[2]{Heures décimales}
Écrire en heures et minutes : $1{,}5$ h ; $2{,}25$ h ; $0{,}1$ h ; $3{,}75$ h.
\end{exercice}\end{blocLivre}
\begin{blocLivre}{ex-12-8}\begin{exercice}[2]{Le piège de la virgule}
Paul pense que $1{,}3$ h et 1 h 30 min représentent la même durée. A-t-il raison ? Écrire 1 h 30 min en heures décimales.
\end{exercice}\end{blocLivre}
\begin{blocLivre}{ex-12-9}\begin{exercice}[2]{La matinée au collège}
La frise représente une matinée. Combien de temps Maya passe-t-elle au collège ? Combien de temps dure la récréation ? Combien de temps reste-t-il pour les cours ?
\label{fig-ex-12-9}\figureGeo{dur-9}
\end{exercice}\end{blocLivre}
\begin{blocLivre}{ex-12-10}\begin{exercice}[2]{Années bissextiles}
Une année est bissextile quand son numéro est un multiple de 4 (sauf certaines années de fin de siècle). L’année 2028 sera-t-elle bissextile ? Combien de jours s’écouleront du 1$^{\text{er}}$ janvier 2027 au 1$^{\text{er}}$ janvier 2029 ?
\end{exercice}\end{blocLivre}

\rubrique{Problèmes}
\begin{blocLivre}{ex-12-11}\begin{exercice}[3]{Le train}
Voici les horaires d’un train. Calculer la durée de chaque partie du trajet, puis la durée totale.\begin{center}\renewcommand{\arraystretch}{1.35}\begin{tabular}{c|ccc}\toprule
Gare & Paris & Lyon & Marseille\\\midrule
Horaire & 7 h 52 & 9 h 48 & 11 h 32\\
\bottomrule\end{tabular}\end{center}
\end{exercice}\end{blocLivre}
\begin{blocLivre}{ex-12-12}\begin{exercice}[3]{La course}
Un coureur termine un marathon en 3 h 25 min 40 s. Son ami arrive 12 min 35 s plus tard. Quel est le temps de l’ami ?
\end{exercice}\end{blocLivre}
\begin{blocLivre}{ex-12-13}\begin{exercice}[3]{Le bus de ville}
Sur une ligne de bus, chacun des 12 tronçons se parcourt en 2 min 45 s. Combien de temps dure le trajet complet ?
\end{exercice}\end{blocLivre}
\begin{blocLivre}{ex-12-14}\begin{exercice}[3]{Les calendriers}
L’année du calendrier julien dure $365{,}25$ jours, soit environ 11 minutes de plus que l’année des saisons. En combien d’années environ ce décalage atteint-il un jour ? Entre 325 et 1582, quel décalage s’est accumulé ? Pourquoi a-t-on supprimé 10 jours en octobre 1582, lors du passage au calendrier grégorien ?
\end{exercice}\end{blocLivre}
\begin{blocLivre}{ex-12-15}\begin{exercice}[3]{Le salaire de la journée}
Un ouvrier travaille 7 h 45 min par jour. Il est payé 14 € de l’heure. Combien gagne-t-il par jour ?
\end{exercice}\end{blocLivre}

\clearpage
